\chapter*{ESQUEMA DEL ÍNDICE GENERAL DE LA TESIS}
\addcontentsline{toc}{chapter}{ESQUEMA DEL ÍNDICE GENERAL DE LA TESIS}
\chaptermark{ESQUEMA DEL ÍNDICE GENERAL}

{\small
\noindent\textbf{CAPÍTULO I: INTRODUCCIÓN Y PLANTEAMIENTO DEL PROBLEMA}\\
1.1. Realidad Problemática\\
1.1.1. La adopción acelerada de la IA Generativa en el desarrollo de software corporativo.\\
1.1.2. La complejidad del dominio del negocio en el sector de Seguros y Reaseguros.\\
1.1.3. La degradación silenciosa de la arquitectura de software y el riesgo operativo resultante.\\
1.2. Formulación del Problema (General y Específicos)\\
1.3. Objetivos de la Investigación (General y Específicos)\\
1.4. Justificación e Importancia (Práctica, Académica, Metodológica e Institucional UNI)\\
1.5. Delimitaciones y Alcance de la Investigación\\[6pt]

\noindent\textbf{CAPÍTULO II: MARCO TEÓRICO Y ESTADO DEL ARTE}\\
2.1. Antecedentes de la Investigación (Internacionales y Nacionales)\\
2.2. Bases Teóricas\\
2.2.1. Deuda Técnica y Mantenibilidad del Software (Modelos SQALE, ISO/IEC 25010).\\
2.2.2. Inteligencia Artificial Generativa y Grandes Modelos de Lenguaje (LLMs).\\
2.2.3. Antipatrones de Código y Code Smells específicos por IAG.\\
2.2.4. Arquitectura de Software para Sistemas de Seguros y Reaseguros.\\
2.3. Definición de Términos Básicos\\[6pt]

\noindent\textbf{CAPÍTULO III: HIPÓTESIS Y OPERACIONALIZACIÓN DE VARIABLES}\\
3.1. Formulación de Hipótesis (General y Específicas)\\
3.2. Operacionalización de Variables\\
3.3. Matriz de Consistencia Académica\\[6pt]

\noindent\textbf{CAPÍTULO IV: METODOLOGÍA DE LA INVESTIGACIÓN}\\
4.1. Tipo y Nivel de Investigación\\
4.2. Diseño de la Investigación\\
4.3. Población, Muestra y Muestreo\\
4.4. Técnicas e Instrumentos de Recolección de Datos\\
4.5. Construcción del Modelo Matemático: Índice de Deuda Técnica Generativa (IDTG)\\
4.6. Procedimiento Experimental y Procesamiento Estadístico\\[6pt]

\noindent\textbf{CAPÍTULO V: DESARROLLO DEL CASO DE ESTUDIO Y RESULTADOS}\\
5.1. Caracterización del Caso de Estudio en Seguros y Reaseguros\\
5.2. Resultados del Análisis Estático de Código\\
5.3. Resultados de las Pruebas Dinámicas de Mantenibilidad\\
5.4. Prueba y Contrastación de Hipótesis\\
5.5. Discusión de Resultados\\[6pt]

\noindent\textbf{CAPÍTULO VI: PROPUESTA DE GOBERNANZA Y GUARDRAILS EN PIPELINES CI/CD}\\
6.1. Arquitectura de Control de Calidad para Aplicaciones Empresariales en Seguros\\
6.2. Reglas de Auditoría Automatizadas y Gatekeepers para Código Generado por IA\\
6.3. Guía de Buenas Prácticas para el Uso de Asistentes de IA en Dominios Financieros\\[6pt]

\noindent\textbf{CONCLUSIONES Y RECOMENDACIONES}\\
\noindent\textbf{REFERENCIAS BIBLIOGRÁFICAS Y ANEXOS}
\par}
